% ============================================================
% ACL 2024 Format – LLM Impersonation Practical Assignment
% Student: Aryan Kathale   |   Teammate impersonated: Ayush
% ============================================================
\documentclass[11pt]{article}
\usepackage[hyperref]{acl}          % ACL 2024 style
\usepackage{times}
\usepackage{latexsym}
\usepackage{amsmath}
\usepackage{graphicx}
\usepackage{booktabs}
\usepackage{multirow}
\usepackage{url}
\usepackage{microtype}
\usepackage{xcolor}

\aclfinalcopy

\title{Impersonating a Teammate Using Large Language Models:\\
       A Practical Study on Persona Mimicry and Evaluation}

\author{
  Aryan Kathale \\
  \textit{[Your Institution]} \\
  \texttt{[your-email@institution.edu]}
}

\begin{document}
\maketitle

% ============================================================
\begin{abstract}
Large language models (LLMs) demonstrate remarkable fluency, raising
societal concerns about their potential for impersonation—mimicking a real
person's language style, opinions, and personality.
This paper investigates how accurately a general-purpose LLM can impersonate
a known individual (\textit{Ayush}) using three prompting strategies:
zero-shot, few-shot, and retrieval-augmented generation (RAG).
We evaluate generated responses against ground-truth answers using three
complementary metrics: ROUGE-L, BERTScore, and sentence-level cosine
similarity.
Results suggest that [TODO: summarise key finding in 1–2 sentences once
experiments are complete].
\end{abstract}

% ============================================================
\section{Introduction}

The rapid advancement of LLMs such as GPT-4 \citep{openai2023gpt4} has
unlocked powerful capabilities in natural language generation, including
the ability to adopt user-defined personas.
While this flexibility has benign applications in education, entertainment,
and accessibility, it also creates serious risks: malicious actors could
use LLMs to impersonate real people, enabling phishing, disinformation,
and identity fraud \citep{zhang2023llmsafety}.

This practical assignment investigates LLM-based impersonation in a
controlled setting.
The subject of impersonation is a consenting teammate, \textbf{Ayush}, who
provided a dataset of 40 personal Q\&A pairs.
The key research questions are:

\begin{enumerate}
  \item How accurately can an LLM reproduce Ayush's language style and
        opinions using different prompting strategies?
  \item Which prompting strategy (zero-shot, few-shot, RAG) yields the
        best impersonation quality?
  \item What do automated evaluation metrics reveal about the strengths
        and limitations of each method?
\end{enumerate}

% ============================================================
\section{Dataset \& Methods}

\subsection{Dataset Creation}
Ayush manually answered 40 questions spanning five categories:
\textit{Daily Life \& Personality} (Q1--10),
\textit{Interests \& Lifestyle} (Q11--18),
\textit{Media \& Preferences} (Q19--25),
\textit{Opinions \& Beliefs} (Q26--34), and
\textit{Deeper / Reflective} (Q35--40).
All responses were written by Ayush without LLM assistance, capturing
authentic personal style.

\paragraph{Data Splits.}
The 40 pairs are divided into three splits:
\begin{itemize}
  \item \textbf{Training (25 pairs, Q1--Q25):} used to build prompts and
        the RAG index.
  \item \textbf{Validation (5 pairs, Q26--Q30):} used to tune prompt
        templates during development.
  \item \textbf{Test (10 pairs, Q31--Q40):} held-out set, never seen
        during prompt design; used for final evaluation only.
\end{itemize}

\subsection{Impersonation Methods}

\paragraph{Method 1 – Zero-shot Prompting.}
A detailed persona description of Ayush (derived from the training data)
is provided as the system prompt.
No example Q\&A pairs are shown.
This tests whether a rich textual persona alone is sufficient for
accurate impersonation.

\paragraph{Method 2 – Few-shot Prompting.}
Five representative training Q\&A pairs (covering all major categories)
are appended to the persona description as in-context examples.
This provides the LLM with concrete stylistic signals.

\paragraph{Method 3 – Retrieval-Augmented Generation (RAG).}
All 25 training Q\&A pairs are embedded using
\texttt{all-MiniLM-L6-v2} \citep{reimers2019sentencebert} and stored in
a FAISS \citep{johnson2019faiss} index.
For each test question, the top-3 semantically similar training pairs are
retrieved and injected as dynamic context.
This method adapts its examples to each question rather than using a
fixed set.

\subsection{Evaluation Metrics}

\paragraph{ROUGE-L \citep{lin2004rouge}.}
Measures the longest common subsequence (LCS) between generated and
reference responses.
ROUGE-L F1 balances precision and recall of shared n-grams.

\paragraph{BERTScore \citep{zhang2020bertscore}.}
Computes token-level cosine similarities between contextual embeddings
(using \texttt{microsoft/deberta-xlarge-mnli}) of generated and reference
text.
Reports precision, recall, and F1.

\paragraph{Cosine Similarity.}
Both generated and reference responses are embedded with
\texttt{all-MiniLM-L6-v2}; their cosine similarity captures overall
semantic alignment beyond surface-level word overlap.

% ============================================================
\section{Results \& Discussion}

Table~\ref{tab:results} compares the three methods across all evaluation
metrics on the 10-item test set.

\begin{table}[h]
\centering
\small
\begin{tabular}{lccc}
\toprule
\textbf{Method} & \textbf{ROUGE-L} & \textbf{BERTScore F1} & \textbf{CosSim} \\
\midrule
Zero-shot  & -- & -- & -- \\
Few-shot   & -- & -- & -- \\
RAG        & -- & -- & -- \\
\bottomrule
\end{tabular}
\caption{Mean scores on the 10-item test set.
         Replace ``--'' with actual values after running experiments.}
\label{tab:results}
\end{table}

\begin{figure}[h]
  \centering
  \includegraphics[width=0.9\columnwidth]{../results/comparison_bar_chart.png}
  \caption{Bar chart comparing impersonation methods across metrics.}
  \label{fig:bar_chart}
\end{figure}

\paragraph{Analysis.}
[TODO: Replace this paragraph with your analysis.]
Discuss (i) which method performed best and why, (ii) whether higher
ROUGE-L correlates with higher semantic similarity, and (iii) categories
or question types where impersonation quality is highest/lowest.

% ============================================================
\section{Ethical Considerations}

LLM-based impersonation poses serious ethical and societal risks.

\paragraph{Privacy.}
Even when the subject consents (as in this study), the resulting model
artefacts could be misused to generate further fake content.
Training data containing personal opinions and lifestyle details should
be treated as sensitive.

\paragraph{Potential Misuse.}
The same techniques used here can be weaponised for spear-phishing,
social engineering, and reputational attacks.
Attackers could impersonate individuals to deceive their contacts or to
fabricate statements.

\paragraph{Consent and Transparency.}
All data used in this study was collected with Ayush's explicit consent.
We recommend that any real-world deployment of impersonation technology
require affirmative consent and clear disclosure to end-users.

\paragraph{Safeguards.}
Potential safeguards include watermarking AI-generated text, building
impersonation detectors, and establishing legal frameworks that criminalise
non-consensual AI impersonation.

% ============================================================
\section{AI Use Declaration}
This report was authored by Aryan Kathale.
The following AI tools were used during the project:

\begin{itemize}
  \item \textbf{GitHub Copilot / GPT-based code completion} was used to
        assist in writing boilerplate Python code (data loading, API
        wrappers, evaluation loops).
        All generated code was reviewed, tested, and adapted by the author.
  \item \textbf{No AI tool was used to write or paraphrase any section of
        this report.}
\end{itemize}

% ============================================================
\bibliography{references}

% ============================================================
\appendix

\section{Full Question List}
\label{app:questions}

All 40 questions used for dataset creation are listed below.
The split assignment is shown in parentheses.

\paragraph{Daily Life \& Personality (Train, Q1--10)}
\begin{enumerate}
  \item What does a typical day look like for you?
  \item What's your morning routine like?
  \item What do you usually do before going to sleep?
  \item How do you deal with stress or pressure?
  \item What's one habit you're proud of?
  \item What's one habit you want to change?
  \item Are you more productive in the morning or at night? Why?
  \item What motivates you to keep going when things get hard?
  \item What kind of environment do you work best in?
  \item How do you usually spend your weekends?
\end{enumerate}

\paragraph{Interests \& Lifestyle (Train, Q11--18)}
\begin{enumerate}\setcounter{enumi}{10}
  \item What are your favorite games or hobbies right now?
  \item Do you prefer staying in or going out? Why?
  \item What's something new you've tried recently?
  \item If you had a completely free week, what would you do?
  \item Do you enjoy traveling? Why or why not?
  \item What's your favorite way to relax?
  \item Are you into fitness or sports? Explain.
  \item What kind of content do you consume the most (YouTube, reels, etc.)?
\end{enumerate}

\paragraph{Media \& Preferences (Train, Q19--25)}
\begin{enumerate}\setcounter{enumi}{18}
  \item What's the last movie or show you watched? Thoughts?
  \item Do you prefer movies, series, or YouTube? Why?
  \item What type of music do you usually listen to?
  \item Do you rewatch shows or prefer new ones?
  \item Who's your favorite creator or influencer right now?
  \item What's a movie or show you think is overrated?
  \item What's one piece of content that really stuck with you?
\end{enumerate}

\paragraph{Opinions \& Beliefs — Validation (Q26--30)}
\begin{enumerate}\setcounter{enumi}{25}
  \item Do you think AI will replace most jobs? Why or why not?
  \item What are your thoughts on social media today?
  \item Is college education still worth it? Explain your view.
  \item Do you think people rely too much on technology?
  \item Privacy vs security---what matters more to you?
\end{enumerate}

\paragraph{Opinions \& Beliefs + Reflective — Test (Q31--40)}
\begin{enumerate}\setcounter{enumi}{30}
  \item What's your opinion on climate change?
  \item How do you feel about AI being used in daily life?
  \item Do you think success depends more on luck or hard work?
  \item What's one opinion you have that most people might disagree with?
  \item What's a goal you're currently working toward?
  \item What kind of person do you want to become in the next few years?
  \item What's something you've learned recently about yourself?
  \item What motivates you in life overall?
  \item How would your close friends describe you?
  \item What impression do you think you give to new people?
\end{enumerate}

\section{Additional Experiment Details}
\label{app:details}

\paragraph{LLM Configuration.}
All experiments use OpenAI's \texttt{gpt-3.5-turbo} (temperature = 0.7,
max\_tokens = 256).
Swapping to \texttt{gpt-4} is supported via the \texttt{--model} flag.

\paragraph{Few-shot Example Selection.}
The five training pairs used as in-context examples are:
Q1 (daily life), Q8 (motivation), Q11 (hobbies), Q17 (fitness), Q20
(media preference).
These were chosen to maximise category coverage.

\paragraph{RAG Index.}
The FAISS index uses \texttt{IndexFlatIP} (inner-product similarity on
L2-normalised embeddings, equivalent to cosine similarity).
The top-3 training pairs are retrieved per query.

\end{document}
